\section{The Access Decision}
\label{sec:access}

\subsection{One question, asked once}

Every endpoint asks the same question of the guard: may this path be read?  The
answer is defined here and every operation is then compared with it.

\begin{definition}[Access]
\label{def:ok}
Let $c=\Canon(\Clean(p))$ be the request path made absolute, cleaned, and
reduced to one alias spelling, and let $\rho=\Real(c)$ be the real location of
what is actually opened.  The type $t$ is the type of the opened object.  Then
\[
  \Ok(p)\;\iff\;
  \InRoots(\rho)\ \land\ \neg\Match(\Deny,\rho,t)\ \land\ \neg\Match(\Deny,c,t)
  \ \land\ \Elig(o),
\]
where $o$ is the opened object and $\Elig(o)$ is the hard-link condition: it holds
unless $o$ is a regular file with $n>1$ names, in which case it holds iff the
index accounts for all of them \emph{now}.  That means that the index is ready,
that it found $n$ distinct names (by folder identity and folded name), that no
deny rule matches the real path of any of them, and that each is still a real
path to $o$ (same device and inode, no symbolic link along it).  If the index is
not ready, or any part cannot be established, $\Elig(o)$ is false: unknown means
refused.  Identity is (device, inode); no content is read.
\end{definition}

Without $\Elig$ the formula would be a path-and-root predicate only.  An allowed
\code{ordinary.txt} that has a second name \code{.env} satisfies the first three
conjuncts and is nevertheless refused, which is the point of the fourth
(Law~\ref{law:linked}).

The check is made on the \emph{real} location (so a symbolic link cannot lead
into a denied place), and also on the \emph{spelled} location (so a link
\emph{named} like a denied file is refused even though its target is fine).  It
is therefore conservative: a path is refused if either name is denied.

\begin{definition}[Open, list, search]
$\Open(p)$ returns the opened object, or $\NotFound$, exactly when $\Ok(p)$
fails or the path does not exist, and otherwise a permission or type error that
mentions nothing about denied paths.  $\List(d)$ is defined when $\Open(d)$ is;
it contains an entry $e$ of $d$ exactly when $\Visible(d,e)$:
\begin{align*}
  \Visible(d,e)\iff{}& \neg\Match(\Deny,\rho_d\cdot e,t)\ \land\
      \neg\MatchBelow(\Hide,\rho_d,\rho_d\cdot e,t)\\
  &\land\ (e\text{ a symbolic link}\Rightarrow\Ok\text{-target}(e))\ \land\
      \Elig_\tau(e),
\end{align*}
where $\Elig_\tau$ is $\Elig$ with one difference: within a single listing or
search a verdict about a file may be reused for at most $\tau=1$ second, after
which the names are checked again.  A listing is therefore a snapshot whose
hard-link verdicts can be up to $\tau$ old; $\Open$ always evaluates $\Elig$
afresh,
and $\rho_d=\Real(d)$: children are judged by where the directory really is,
not by the link it was reached through.  $\Search(d,q)$ walks breadth-first from
$d$, entering only visible unlinked directories whose real location is still the
name they were queued under, and reports the visible entries whose folded name
contains $\fold{q}$.  It also reports whether it was \emph{complete}: not if a
folder that was allowed could not be opened or read to the end, or was replaced by a
link before it was entered; denied, vanished and cloud-only folders do not count.
\end{definition}

\subsection{Laws}

\begin{law}[Denied is absent]
\label{law:absent}
If $\neg\Ok(p)$ because of a deny rule, then $p$ is not opened, not listed by
$\List(\Parent(p))$, and not reported by any $\Search$.
\Tests{guard/laws_test.go::TestLawDeniedIsAbsentEverywhere; guard/fuzz_test.go (independent inode and marker oracle)}
\end{law}

\begin{law}[Every spelling the OS resolves is refused]
\label{law:differential}
For every secret $s$ and every spelling $p$ such that the operating system
resolves $p$ to the same object as $s$ (case, folding aliases, decomposed accents,
doubled separators, dot segments, trailing separators, the data-volume and
symbolic-link aliases, named forks), every operation refuses $p$.  The premise is
decided by the OS (same inode), not by \texttt{fsb}, so the law does not depend on
\texttt{fsb}'s own idea of a name.  Dually, every such spelling of an ordinary file
inside the root is served.
\Tests{guard/differential\_test.go (both directions; fails 1,536 times if the folding fix is reverted)}
\end{law}

\begin{law}[Every name of a denied file is denied]
\label{law:linked}
If a regular file has several names and one of them is matched by a deny rule,
then $\Open$ refuses the file under every name, and no listing shows it.  If the
names cannot be determined, or the index does not account for every one of them
(as many distinct names as the file has links, each still a real path to it), the
file is refused.  Identity is the pair
(device, inode), never a hash of the contents, so an ordinary copy of a secret is
not treated as the secret and no file is read to decide.
\Tests{guard/hardlink\_test.go (all four rule kinds, links made after the index was built, links to files outside the roots, refusal while the index builds, bounded walk, names renamed or moved, names outside every location, the retry walk and cooldown, the one-second verdict reuse in listings and searches); e2e ``a hard link to a denied file is refused under any name''}
\end{law}

\begin{law}[Indistinguishable from missing]
\label{law:indist}
For every endpoint $f$, a denied path $p$ and a missing path $q$ receive the
same status and body: $f(p)=f(q)$.  This includes unreadable files inside a
denied folder, which must not answer 403.
\Tests{guard/laws_test.go::TestLawDeniedIsAbsentEverywhere; server/m2_test.go; e2e ``answers 404, with no content, for denied paths on every endpoint''}
\end{law}

\begin{law}[Listed implies openable]
\label{law:listed-open}
If $e$ is listed in $\List(d)$ and is not a dangling link, then $\Open(d\cdot e)$
succeeds, \emph{provided the file system does not change between the two
operations and the operating system's own access and type conditions hold}
(permissions, a regular file or folder).  This is a snapshot property: within
$\tau$ of a change, a listing may still show a multi-name file that a fresh
$\Open$ refuses (the \emph{Hard links} boundary).  The converse is false by
design: hidden entries are not listed but can be opened.
\Tests{guard/laws_test.go::TestLawListedImpliesOpenable}
\end{law}

\begin{law}[Found implies listed]
\label{law:found-listed}
If $\Search$ reports $m$, then $m$ is an entry of $\List(\Parent(m))$, for a
file system that does not change between the two observations (and, for
multi-name files, to within $\tau$).  Search therefore cannot show what browsing
would not, because both go through one function.
\Tests{guard/laws_test.go::TestLawSearchResultsAreListed}
\end{law}

\begin{law}[Hidden is reachable]
\label{law:hidden-reachable}
If $\Match(\Hide,p,t)$ but $\Ok(p)$ holds, then $\Open(p)$ succeeds.  Hide rules
are cosmetic and never restrict direct access.
\Tests{guard/guard_test.go (hidden entries); e2e ``lists the home folder without denied or ignored entries''}
\end{law}

\begin{law}[Containment]
\label{law:containment}
$\Ok(p)\Rightarrow\exists r\in\Roots:\ \Real(p)\text{ is }r\text{ or lies beneath }r$,
where ``is $r$ or lies beneath $r$'' is the identity relation $\InRoots$ of
Law~\ref{law:root-identity}, not a comparison of spellings.  (Folding belongs to
the name and rule model, where over-matching is safe; it is not used for this
\emph{allow} decision.)
\Tests{guard/guard_test.go; guard/firmlink_darwin_test.go}
\end{law}

\begin{law}[Containment is by identity]
\label{law:root-identity}
$\InRoots(\rho)$ holds iff the leading components of $\rho$ that correspond to a
root name \emph{the same directory} as that root (same device and inode).
Comparing spellings would be right on a case-insensitive volume, but on a
case-sensitive one it would admit a sibling folder that differs from the root
only in case.  This is an \emph{allow} decision, so unlike the deny rules it must
not over-match.
\Tests{guard/casesens\_test.go (needs FSB\_CASE\_SENSITIVE\_VOLUME); guard/guard\_test.go::TestInRoots}
\end{law}

\begin{law}[Alias reduction is by name equality]
\label{law:alias-fold}
$\Canon$ recognizes the data-volume prefix and the firmlink directories
component by component, comparing components with $\Fold$, so no spelling of
\code{/System/Volumes/Data/Users} escapes reduction.  Slicing by byte length
does not have this property, because a folded spelling can be longer or shorter.
\Tests{guard/foldalias\_test.go::TestRootOfSlashHonoursDenyRulesThroughFoldedFirmlinkSpellings; guard/firmlink\_darwin\_test.go}
\end{law}

\begin{law}[The session is scoped to a random path]
\label{law:cookie-path}
Browsers scope a cookie by host and never by port, so a cookie for
\code{127.0.0.1} would accompany every request to every other web server on that
address.  All URLs of \texttt{fsb} therefore live under a random 128-bit prefix
$/\pi/$ fixed at launch; the cookie is set with \code{Path=/$\pi$/}; and a
request whose path is not under the prefix is refused exactly like an
unauthenticated one, whatever cookie it carries.  A cookie that is sent only for
paths under $/\pi/$ is never sent to another server, which has no such path.
The frontend forms every request relative to the page's own directory.
\Tests{httpguard/httpguard\_test.go::TestOnlyURLsUnderThePrefixExist, TestSessionCookieIsScopedToThePrefix; e2e ``the session cookie is scoped to fsb's own prefix'' (fails if the path is widened); web/src/lib/laws.test.ts::apiBase}
\end{law}

\begin{law}[Canonical form is stable]
\label{law:canon}
$\Canon(\Canon(p))=\Canon(p)$.
\Tests{guard/laws_test.go::TestLawCanonPathIsIdempotent}
\end{law}

\begin{law}[Bounded work]
\label{law:bounded}
$\Search$ examines at most $V$ directory entries (excluded ones included) and
reports at most $M$ results,
and reports truncation when a bound stopped it; it never follows symbolic links
to directories, so it terminates on cyclic links.
\Tests{guard/search_test.go; e2e (100,000-file directory, optional)}
\end{law}

\subsection{Boundaries}

\begin{boundary}[Time of check]
The file is opened first and its real path is then resolved and confirmed to
name the same object (\code{SameFile}); the decision is taken on the object that
is actually served.  A rename between opening and resolving is treated as a
race and refused.
\end{boundary}

\begin{boundary}[Response timing]
A denied path can be answered slightly faster than a missing one.  Status and
body are identical; timing is not equalized.
\end{boundary}

\begin{boundary}[Other processes]
Anything that runs as the user can read the same files.  \texttt{fsb} is not a
sandbox against local malware.
\end{boundary}
